\clearpage

% 5. AREAS
\section{Áreas de Aprendizado}

Este capítulo apresenta as áreas matemáticas centrais que guiarão as atividades do clube,
selecionadas por sua relevância direta para o cursus da 42 e para o desenvolvimento do
raciocínio lógico-matemático dos cadetes.

\subsection{Área 1 --- Lógica e Raciocínio Matemático}

A lógica formal é a espinha dorsal da programação e da matemática. Esta área cobre
proposições, conectivos lógicos, tabelas-verdade, quantificadores, implicação, equivalência
e demonstrações. Dominar lógica permite ao cadete entender o funcionamento interno dos
programas e construir argumentos rigorosos.

\subsection{Área 2 --- Matemática Discreta}

A matemática discreta é o campo mais diretamente aplicável à ciência da computação. Inclui
teoria ingênua dos conjuntos, relações, funções, combinatória, teoria dos grafos e indução
matemática. Esses tópicos aparecem constantemente em algoritmos, estruturas de dados e
análise de complexidade.

\subsection{Área 3 --- Álgebra e Aritmética}

Uma base sólida em álgebra é fundamental para qualquer área da matemática aplicada. Esta
área cobre manipulação algébrica, teoria dos números, aritmética modular, divisibilidade,
MDC, MMC e sistemas de equações.

\subsection{Área 4 --- Probabilidade e Estatística}

Probabilidade e estatística são cada vez mais relevantes na computação moderna,
especialmente em algoritmos aleatorizados, análise de dados e inteligência artificial.
Esta área abrange probabilidade básica, distribuições, esperança, variância e análise
descritiva.

\subsection{Área 5 --- Algoritmos e Complexidade}

Esta área conecta a matemática diretamente à programação. Abrange análise de algoritmos,
notação Big-$\mathcal{O}$, complexidade de tempo e espaço, recorrências e técnicas de
design de algoritmos (divisão e conquista, programação dinâmica, algoritmos gulosos).
